% DO NOT REMOVE SPACES BETWEEN LINES!

\begin{center}\includegraphics[width=0.6\textwidth]{liu_logo.png}\end{center}

\vspace{-1mm}
\begin{center}\Large{\textbf{\color{bluePoli}{\title}}}\end{center}

\vspace{-0.2cm}
\begin{center}\noindent\fontsize{0.4cm}{0.5cm}\selectfont \bfseries \textsc{\color{bluePoli} Report for Design Task in \course}\end{center}

\vspace{-0.2cm}
\begin{center} \large{\textbf{Authors:}} \end{center}
\begin{center}\noindent\!\!\!\large{\textbf{\author}}\end{center}
\vspace{0.05cm}

\small \normalfont

\vspace{11pt}

\centerline{\rule{1.0\textwidth}{0.4pt}}

\begin{center}
\begin{minipage}[t]{.24\textwidth}
\begin{minipage}{.9\textwidth}
\noindent
\scriptsize{\textbf{Professor:}} \\
\advisor \\
\\
\textbf{Assistant Professor:} \\ 
\firstcoadvisor \\
\secondcoadvisor \\ 
\textbf{Academic Year:} \\
\YEAR \\
\\
\end{minipage}
\end{minipage}% This must go next to `\end{minipage}`
\begin{minipage}{.74\textwidth}
\noindent \textbf{\color{bluePoli} Abstract:} {
    The following report consists of the design and mathematical process which we underwent to design the hydraluic circuit for a lorry crane. The system contains two cylinder, and it was our task to determine the operating conditions for the cylinders and pump, and pick-out the appropriate components for the task. The dimensioning also required us to consider a work-case scenario, which informed our decision during part-selection.
    }
\end{minipage}
\end{center}

\vspace{12pt}